\documentclass{article}
\usepackage{amsmath, amssymb, amsthm}
\usepackage{hyperref}
\usepackage{geometry}
\geometry{margin=1in}

\title{Erd\H{o}s Problem \#1108}
\date{Last updated: 2025-11-17}
\author{Source: erdosproblems.com}

\begin{document}
\maketitle

\noindent\textbf{Status:} OPEN \quad \textbf{No prize}

\noindent\textbf{Tags:} number theory, factorials

\medskip
\noindent\textbf{Problem Statement:}

Let\[A = \left\{ \sum_{n\in S}n! : S\subset \mathbb{N}\textrm{ finite}\right\}.\]If $k\geq 2$, then does $A$ contain only finitely many $k$th powers? Does it contain only finitely many powerful numbers?

\bigskip
\noindent\textbf{Remarks:}

Asked by Erd\H{o}s at Oberwolfach in 1988. It is open even whether there are infinitely many squares of the form $1+n!$ (see [398]).

This was motivated in part by a problem of Mahler which he discussed with Erd\H{o}s a few days before his death in 1988: if $k\geq 5$ and\[A_k= \left\{ \sum_{n\in S}k^n : S\subset \mathbb{N}\textrm{ finite}\right\}\]then does $A_k$ contain only finitely many squares? Mahler showed that there are infinitely many squares in $A_k$ for $k\leq 4$, and found only one square for $k\geq 5$, namely\[1+7+7^2+7^3=400.\]Brindza and Erd\H{o}s \cite{BrEr91} proved that, for any $r$, if $n_1!+\cdots+n_r!$ is powerful then $n_1\ll_r 1$.

\bigskip
\noindent\textbf{References:}

[BrEr91] Brindza, B. and Erd\H{o}s, P., On some {D}iophantine problems involving powers and
factorials. J. Austral. Math. Soc. Ser. A (1991), 1--7.

\bigskip
\noindent\small{Source: \url{https://www.erdosproblems.com/1108}}
\end{document}
